\begin{afterword}
\summaryhead

The post explains why a sequence on reductionism includes posts about joy in the merely real. Taking things apart can feel like removing magic from the world, and the author wants to leave readers a line of retreat. The posts also oppose ``Hollywood Rationality,'' the view that understanding the rainbow subtracts its beauty.

The deeper theme is ``binding yourself to reality.'' Hope spent on the lottery should not be given up but moved to a job, a school or a business, since the lottery is not hard but ``un-actionable.'' In the same way, emotional energy attached to magic or to bad explanations of rainbows should be redirected into reality, not destroyed. Emotions that follow from one's best picture of the truth are not irrational, as ``Feeling Rational'' argued. If you want to fly, give up flying potions and build an airplane. The post ends: ``Science puts the fun back into life.''

\responsehead

The post is short and states an attitude, and it should be judged on the few claims it makes. It does what it sets out to do: it tells the reader why posts about feeling appear in a sequence about reductionism, and it gives the attitude a name.

Its central idea, that feeling can be moved from impossible objects to real ones, is offered as advice in a metaphor, ``binding it into the realms of reality,'' not as a finding.

The examples do not test it. The lottery example restates an earlier post, whose own evidence is discussed in the notes on that post. The airplane is the easy case: it gives exactly what the potion promised, flight, so nothing has to be transferred between different goods. The case the post is about is the rainbow, where what replaces the magic is a different pleasure, the physics of light. Whether the feeling that went to the magic will go to the physics is the question, and the post does not show that it will.

The reply to the rainbow's critics is not new, and the post does not claim that it is. It is the thesis of Richard Dawkins's \textsc{Unweaving the Rainbow} (1998), named after the same lines of Keats, and ``Joy in the Merely Real,'' two days earlier, takes an epigraph from Feynman to the same effect.

In short: a clear statement of an attitude, offered as advice, whose main example is the easy case, in which the real thing gives what the magic promised.
\end{afterword}
